\chapter{CI/CD with Databricks Asset Bundles, Terraform, and Cost Optimization}
\index{Asset Bundles}\index{databricks.yml}\index{Terraform}\index{CI/CD}\index{GitHub Actions}%
\index{infrastructure as code}\index{dev staging prod}\index{cost optimization}\index{system.billing.usage}%
\begin{objectives}
\begin{itemize}
    \item Define jobs, pipelines, and compute as code with Databricks Asset Bundles.
    \item Structure bundle targets (dev/staging/prod) with variables and mode constraints.
    \item Provision workspace resources with the Terraform Databricks provider and remote state.
    \item Build a GitHub Actions pipeline: PR validation, staging gate, prod promotion.
    \item Produce per-domain chargeback from \texttt{system.billing.usage} and act on it.
    \item Design a delivery and FinOps operating model for a multi-team platform.
\end{itemize}
\end{objectives}

\begin{crosscloud}
Asset Bundles are Databricks' \textbf{native IaC}; Terraform is the cross-cloud constant; the
CI/CD pattern is identical everywhere.

\vspace{2mm}
{\footnotesize
\begin{tabular}{@{}p{3.0cm}p{3.4cm}p{3.2cm}p{3.2cm}@{}}
\toprule
\textbf{Capability} & \textbf{Databricks} & \textbf{AWS} & \textbf{Azure / GCP / Fabric} \\
\midrule
Native IaC & Asset Bundles (\texttt{databricks.yml}) & CloudFormation / CDK & Bicep / Deployment Manager / Fabric Git + Terraform \\
Cross-platform IaC & Terraform provider & Terraform & Terraform \\
Validate-before-deploy & \texttt{bundle validate} & cfn-lint / CDK synth & \texttt{terraform plan} / what-if \\
CI auth & Service principal + OAuth M2M & OIDC $\rightarrow$ IAM role & OIDC $\rightarrow$ workload identity \\
State & Terraform remote state & S3 + DynamoDB lock & Storage account / GCS backend \\
Cost telemetry & \texttt{system.billing.usage} & CUR 2.0 & Cost Management / billing export / Capacity Metrics \\
\bottomrule
\end{tabular}}

\vspace{1mm}
\textbf{Snowflake}'s story is Terraform-first plus database-change tools (schemachange, dbt);
the lesson that transfers everywhere: \emph{the merge is the deploy, and the pipeline is the only
writer to prod}.
\end{crosscloud}

\section{Week 23 Roadmap and Mental Model}
Twenty-two weeks of platform skills converge on the delivery question: how does a change travel from a laptop to production \emph{safely}? Week 23's answer is infrastructure as code (Asset Bundles for Databricks resources, Terraform for the workspace perimeter) behind a CI/CD pipeline with a staging gate---plus the FinOps loop that keeps the whole estate affordable \cite{databricksAssetBundles,terraformDatabricks}.

The mental model: \textbf{production is a build artifact}. Nobody clicks in prod; a merge to \texttt{main} triggers validation, staging deployment, tests, and---only then---prod deployment. The bundle is the unit of deployment; the target is the environment; the pipeline is the policy.

\begin{definitionbox}
\textbf{A Databricks Asset Bundle} is a versioned declaration (\texttt{databricks.yml} plus resource files) of jobs, pipelines, and other workspace objects, deployed to named \textbf{targets} (dev/staging/prod) with variable substitution, deployment state tracking, and \texttt{validate}/\texttt{deploy}/\texttt{run} lifecycle commands.
\end{definitionbox}

\begin{keyterms}
\textbf{Target}: a named deployment environment with its own workspace, variables, and mode. \textbf{Mode}: \texttt{development} (prefixed, paused schedules) or \texttt{production} (unpaused, locked). \textbf{Variable}: a bundle input substituted per target. \textbf{Remote state}: Terraform's shared, locked state backend. \textbf{Chargeback}: attributing platform cost to domains from billing system tables.
\end{keyterms}

\section{Monday: Asset Bundles}
\subsection{Anatomy}
A bundle declares resources (jobs, pipelines, dashboards, models) and targets; \texttt{databricks bundle validate} checks the declaration against the workspace schema before anything deploys; \texttt{bundle deploy} syncs code and applies resources; \texttt{bundle run} executes a resource in a target \cite{databricksAssetBundles}.

\begin{lstlisting}[language=json,caption={The course bundle: one job, one pipeline, three targets.},label={lst:ch23-bundle}]
bundle:
  name: governed-platform

variables:
  catalog: {default: de_training}

targets:
  dev:
    mode: development        # prefixed names, paused schedules
    default: true
    workspace: {host: https://dev-ws.azuredatabricks.net}
  staging:
    mode: production
    workspace: {host: https://stg-ws.azuredatabricks.net}
    variables: {catalog: de_staging}
  prod:
    mode: production
    workspace: {host: https://prd-ws.azuredatabricks.net}
    variables: {catalog: de_prod}
    run_as: {service_principal_name: sp-etl-prod}

resources:
  pipelines:
    medallion:
      name: medallion-${bundle.target}
      catalog: ${var.catalog}
      libraries:
        - notebook: {path: ../src/pipelines/medallion.py}
  jobs:
    nightly_elt:
      name: nightly-elt-${bundle.target}
      tasks:
        - task_key: pipeline
          pipeline_task: {pipeline_id: ${resources.pipelines.medallion.id}}
        - task_key: quality_gate
          depends_on: [{task_key: pipeline}]
          notebook_task: {notebook_path: ../src/quality_gate}
\end{lstlisting}

\subsection{Targets and Modes}
Development mode prefixes resource names and pauses schedules so ten engineers' dev deployments coexist without collisions; production mode unpins schedules and enforces \texttt{run\_as}. The discipline that prevents the classic incident: \emph{prod is deployed only by CI}, from the pinned merge commit, after staging passes.

\begin{pitfall}
\textbf{Deploying prod from a laptop.} Local prod deploys bypass the staging gate, the audit trail, and the pinned commit. Remove human deploy rights to prod targets; CI's service principal holds them. The exception path (break-glass deploy) is documented, credentialed separately, and alarmed.
\end{pitfall}

\section{Tuesday: Terraform for the Workspace Perimeter}
\subsection{What Terraform Owns versus What Bundles Own}
The split this book uses: \textbf{Terraform} provisions the slow-moving perimeter---workspaces, service principals, groups, metastore bindings, cluster policies, secret scopes; \textbf{bundles} deploy the fast-moving application---jobs, pipelines, dashboards. Both are code; both go through the same review; neither is clicked in production \cite{terraformDatabricks}.

\begin{lstlisting}[language=json,caption={Terraform perimeter for the prod target.},label={lst:ch23-terraform}]
terraform {
  backend "azurerm" {                     # remote, locked state
    resource_group_name  = "rg-de-terraform"
    storage_account_name = "stdetfstate"
    container_name       = "tfstate"
    key                  = "databricks-prod.tfstate"
  }
  required_providers {
    databricks = { source = "databricks/databricks" }
  }
}

provider "databricks" { host = var.workspace_host }

resource "databricks_service_principal" "etl_prod" {
  display_name = "sp-etl-prod"
}

resource "databricks_cluster_policy" "training" {
  name       = "de-training-policy"
  definition = jsonencode({
    autotermination_minutes = { type = "fixed", value = 30 }
    num_workers             = { type = "range", maxValue = 4 }
  })
}

resource "databricks_grants" "gold_readers" {
  schema = "de_prod.gold"
  grant { principal  = "analysts"
          privileges = ["USE_SCHEMA", "SELECT"] }
}
\end{lstlisting}

\subsection{State Discipline}
Terraform state is the record of what exists: remote (never a laptop file), locked (no concurrent applies), and access-controlled (it describes your security estate). A lost or forked state file is how ``Terraform wants to delete the metastore binding'' happens.

\section{Wednesday: The Pipeline and the FinOps Loop}
\subsection{GitHub Actions Shape}
\begin{lstlisting}[language=json,caption={CI/CD: validate on PR, deploy to staging on merge, gate prod.},label={lst:ch23-actions}]
name: platform-cicd
on:
  pull_request: {branches: [main]}
  push: {branches: [main]}
env:
  DATABRICKS_CLIENT_ID: ${{ vars.SP_CLIENT_ID }}
  DATABRICKS_CLIENT_SECRET: ${{ secrets.SP_CLIENT_SECRET }}
jobs:
  validate:
    runs-on: ubuntu-latest
    steps:
      - uses: actions/checkout@v4
      - uses: databricks/setup-cli@main
      - run: databricks bundle validate -t staging
  deploy_staging:
    needs: validate
    if: github.ref == 'refs/heads/main'
    runs-on: ubuntu-latest
    steps:
      - uses: actions/checkout@v4
      - uses: databricks/setup-cli@main
      - run: |
          databricks bundle deploy -t staging
          databricks bundle run nightly_elt -t staging
          python tests/smoke_staging.py        # contract + row-count checks
  deploy_prod:
    needs: deploy_staging
    environment: production                    # requires reviewer approval
    runs-on: ubuntu-latest
    steps:
      - uses: actions/checkout@v4
      - uses: databricks/setup-cli@main
      - run: databricks bundle deploy -t prod
\end{lstlisting}

\subsection{FinOps from System Tables}
Weekly cost practice: \texttt{system.billing.usage} by SKU and tag, per-domain chargeback, and the two-optimization rule (every monthly review implements at least two cost changes: auto-termination enforcement, serverless conversion, trigger-cadence rightsizing, storage lifecycle) \cite{databricksSystemTables}.

\begin{lstlisting}[language=SQL,caption={Per-domain monthly chargeback.},label={lst:ch23-chargeback}]
SELECT custom_tags.domain,
       sku_name,
       ROUND(SUM(usage_quantity), 1) AS dbus,
       ROUND(SUM(usage_quantity * list_price), 0) AS est_usd
FROM system.billing.usage
JOIN system.billing.list_prices USING (sku_name, cloud, price_start_time)
WHERE usage_date >= date_trunc('month', current_date()) - INTERVAL 1 MONTH
GROUP BY custom_tags.domain, sku_name
ORDER BY est_usd DESC;
\end{lstlisting}

\section{Thursday Hands-on Project: Bundle, Pipeline, and the Cost Report}
Convert the Week 12 job into a bundle with dev/prod targets, drive it from GitHub Actions through a staging gate, provision the perimeter with Terraform, and ship the weekly DBU report.
\index{hands-on lab}\index{Asset Bundles!lab}\index{CI/CD!lab}

\begin{lstlisting}[language=PowerShell,caption={Week 23 lab: from hand-built job to gated deployment.},label={lst:ch23-lab}]
# 1. Scaffold and validate
databricks bundle init   # or author databricks.yml per lst:ch23-bundle
databricks bundle validate -t dev
databricks bundle deploy  -t dev
databricks bundle run nightly_elt -t dev

# 2. Terraform the perimeter (lst:ch23-terraform)
terraform init; terraform plan -out=tfplan; terraform apply tfplan

# 3. Wire .github/workflows/platform-cicd.yml (lst:ch23-actions);
#    PR -> validate; merge -> staging deploy + smoke; approval -> prod.

# 4. The Friday habit, automated:
databricks bundle run weekly_cost_report -t prod
\end{lstlisting}

\subsection*{Deliverables}
\begin{itemize}
    \item The bundle with dev/staging/prod targets and variables; validation output.
    \item The Terraform configuration with remote, locked state.
    \item The Actions run history: PR validation, staging gate, approved prod deploy.
    \item The chargeback report with two implemented optimizations.
\end{itemize}

\subsection*{Acceptance Criteria}
\begin{itemize}
    \item A clean clone deploys the platform to dev with two commands (\texttt{terraform apply}, \texttt{bundle deploy}).
    \item A deliberately broken bundle fails \texttt{validate} in the PR and never reaches staging.
    \item Prod deploys only via the approved environment; the audit trail shows it.
    \item The cost report attributes $\geq$95\% of DBUs to a domain tag.
\end{itemize}

\subsection*{Stretch Goals}
\begin{itemize}
    \item Add the lineage impact check from Chapter 22 as a PR comment step.
    \item Add bundle tests (notebook unit tests) that run in the staging job.
    \item Implement the break-glass prod deploy path with its own alarm.
\end{itemize}

\section{Friday Use Case: Thirty Pipelines, Three Environments, Monthly Accountability}
A platform team ships 30 pipelines for five business domains across dev/staging/prod, with a per-domain monthly budget and a release-train cadence. Design the delivery and FinOps operating model.
\index{operating model}\index{FinOps}\index{release train}

\subsection{Requirements and Assumptions}
Domains ship weekly on the train; urgent fixes ride a documented hotfix path. Finance expects per-domain monthly variance explanations. Two production incidents last quarter were untracked manual changes.

\begin{figure}[H]
    \centering
    \resizebox{\textwidth}{!}{%
    \begin{tikzpicture}[
        box/.style={rectangle, rounded corners, draw=primary, thick, fill=primary!7, align=center, minimum width=2.6cm, minimum height=0.95cm, font=\small\bfseries},
        gate/.style={rectangle, rounded corners, draw=orange!80!black, thick, fill=orange!10, align=center, minimum width=2.4cm, minimum height=0.85cm, font=\scriptsize\bfseries},
        flow/.style={-{Latex[length=2mm]}, draw=primary, thick}
    ]
        \node[box] (pr) at (0,0) {PR\\validate +\\lineage check};
        \node[box] (stg) at (3.6,0) {Staging\\deploy +\\smoke tests};
        \node[gate] (gate) at (7.2,0) {Approval\\(release\\manager)};
        \node[box] (prod) at (10.4,0) {Prod deploy\\(CI only)};
        \node[box] (fin) at (13.6,0) {FinOps\\chargeback +\\variance};
        \draw[flow] (pr) -- (stg);
        \draw[flow] (stg) -- (gate);
        \draw[flow] (gate) -- (prod);
        \draw[flow] (prod) -- (fin);
    \end{tikzpicture}%
    }
    \caption{The release train: every change rides the same validated, gated track; cost closes the loop monthly.}
    \label{fig:ch23-train}
\end{figure}

\subsection{Architecture Decisions}
\textbf{One track, no exceptions.} Every resource change rides PR $\rightarrow$ staging $\rightarrow$ approval $\rightarrow$ prod; the hotfix path is the same track with an expedited approval, so the audit trail survives the emergency. The two manual-change incidents become impossible by removing human prod-deploy rights and alarming on drift (a scheduled job compares deployed state to the bundle registry).

\textbf{Cost as a domain conversation.} Chargeback lands monthly per domain from \cref{lst:ch23-chargeback}; variance $>$15\% requires a note; the two-optimization rule keeps the estate's unit cost falling even as volume grows. Domains see their DBU dashboard---cost awareness without central nagging.

\begin{pitfall}
\textbf{The untracked manual fix.} Every ``quick fix in the UI'' forks production from the code. Detect it (drift job), reverse it (the bundle redeploys over it), and fix the loop time that tempted it---if the train takes two hours, hotfixes will always look faster, and they always cost more than they save.
\end{pitfall}

\subsection{Risks and Mitigations}
\begin{table}[H]
\centering
\small
\begin{tabular}{@{}p{4.6cm}p{7.6cm}@{}}
\toprule
\textbf{Risk} & \textbf{Mitigation} \\
\midrule
Manual prod change under pressure & Humans lack prod deploy rights; drift job alarms on divergence \\
Staging no longer predicts prod & Scheduled staging refresh; parity checks in the pipeline \\
Terraform state corrupted or locked & Remote versioned backend; documented force-unlock procedure \\
Chargeback tags missing on new compute & Tag enforcement in cluster policy, not in documentation \\
Approval becomes rubber-stamping & Two-reviewer rule for schema or grant changes \\
\bottomrule
\end{tabular}
\caption{Week 23 scenario risks and mitigations.}
\label{tab:ch23-risks}
\end{table}

\section{Design Decision Guide}
\index{decision guide!Week 23}%
Delivery decisions are policy made executable; optimize for the audit trail surviving the emergency.

\begin{table}[H]
\centering
\small
\begin{tabular}{@{}p{3.4cm}p{4.4cm}p{4.6cm}@{}}
\toprule
\textbf{Decision} & \textbf{Choose the first when} & \textbf{Choose the second when} \\
\midrule
Bundles vs.\ Terraform & Application resources (jobs, pipelines) & Perimeter (workspaces, principals, policies) \\
Staging gate vs.\ direct prod & Anything user-facing or scheduled & Dev-only sandboxes \\
Service principal in CI vs.\ PAT & Always the principal & Never a personal token in CI \\
Hotfix via pipeline vs.\ manual fix & Always the pipeline, expedited & Never manual---drift is the incident \\
Chargeback vs.\ showback & Budgets have owners & Early-stage awareness building \\
Weekly train vs.\ continuous release & Team coordination matters & Mature trunk-based teams \\
\bottomrule
\end{tabular}
\caption{Week 23 decision guide.}
\label{tab:ch23-decisions}
\end{table}

The two-command redeploy (Terraform apply, bundle deploy) is the health metric of the whole delivery system: when it stops being true, the estate is drifting from its code.

\section{Operational Guidance for Week 23}
\index{operational guidance!Week 23}%
\subsection{Performance and Reliability}
Measure the delivery system itself: deploy frequency, PR-to-prod lead time, validation failure rate, and staging-to-prod divergence time. A staging environment that drifts from prod (data volume, configuration) stops predicting prod---refresh it on a schedule. The drift-detection job (deployed state versus bundle registry) runs daily; any diff is either an untracked change or a failed deploy, and both are incidents-lite.

\subsection{Security and Governance Checklist}
\begin{itemize}
    \item CI service principal holds deploy rights; humans hold approval rights---separate duties.
    \item Terraform state backend access-controlled and versioned.
    \item Bundle variables containing environment secrets resolved from CI secret stores.
    \item Prod targets locked to pipeline deployment; break-glass path alarmed.
\end{itemize}

\subsection{Troubleshooting Playbook}
\begin{table}[H]
\centering
\small
\begin{tabular}{@{}p{3.6cm}p{4.2cm}p{4.8cm}@{}}
\toprule
\textbf{Symptom} & \textbf{Likely cause} & \textbf{First action} \\
\midrule
Validate passes, deploy fails & Permission gap on the target workspace & Deploy as the CI principal; compare grants per target \\
Staging green, prod red & Environment/config drift & Diff target variables and workspace configuration \\
Terraform state locked & Aborted apply left the lock & Force-unlock by lock ID after confirming no live apply \\
Bundle variable missing in one target & Variable declared without default & Add per-target overrides or a safe default \\
\bottomrule
\end{tabular}
\caption{Week 23 troubleshooting quick reference.}
\label{tab:ch23-trouble}
\end{table}

\section{Interview Q\&A}
\subsection{Concepts and Trade-offs}
\begin{enumerate}
    \item \textbf{Q: Asset Bundles versus Terraform---what goes where?}\\
    A: Terraform for the slow perimeter (workspaces, principals, policies, metastores); bundles for the fast application layer (jobs, pipelines, dashboards) with per-target variables and modes. Both in Git, both through the same review.
    \item \textbf{Q: What does \texttt{bundle validate} actually catch?}\\
    A: Schema errors, missing variables, bad resource references, and permission problems---before any cloud change. It is the compile step for your platform.
    \item \textbf{Q: Why remote, locked Terraform state?}\\
    A: State is the record of what exists. Local state forks under team/CI use; a forked state is how plans propose recreating---or deleting---your production estate.
    \item \textbf{Q: How do you keep humans out of prod without slowing hotfixes?}\\
    A: The same pipeline with an expedited approval gate: the hotfix is a PR with a faster review SLA, not a different process. Audit and staging evidence survive the emergency.
    \item \textbf{Q: Where does chargeback data come from, and what breaks it?}\\
    A: \texttt{system.billing.usage} plus tags (cluster policy-enforced, bundle-injected). Untagged compute breaks attribution---which is why tag enforcement lives in policy, not in a wiki.
    \item \textbf{Q: What is the first signal your delivery process is failing?}\\
    A: Drift between deployed state and the bundle registry---caught by the drift job. Manual changes are not a people problem; they are a missing-guardrail problem.
    \item \textbf{Q: How do you attribute a cost spike to a change?}\\
    A: Join the billing spike window to the deployment history: bundle deploys and Terraform applies are logged, so the cost curve diffs against the release curve.
    \item \textbf{Q: How do you convince a team to stop deploying by hand?}\\
    A: Make the pipeline faster than the manual path and prove it once: a two-command redeploy and a five-minute PR-to-staging loop. Teams abandon manual deploys when the governed path is also the easiest one.
\end{enumerate}

\section{Further Reading}
Asset Bundles are documented in \cite{databricksAssetBundles}; the Terraform provider in \cite{terraformDatabricks}; billing system tables in \cite{databricksSystemTables}. The jobs being deployed come from Chapter 12 \cite{databricksWorkflows}; the service principals CI runs as come from Chapter 21 \cite{databricksSecurity}.

\section*{Key Takeaways}
\begin{itemize}
    \item Production is a build artifact: bundles + Terraform in Git, deployed by CI, gated by staging.
    \item Targets and modes make dev collisions and prod schedules structurally impossible to confuse.
    \item Terraform state is remote, locked, and access-controlled---it describes your security estate.
    \item Hotfixes ride the same track faster, never a different track.
    \item FinOps is a monthly loop with per-domain chargeback and two implemented optimizations---enforced by tags in policy.
\end{itemize}

\section{Exercises}
\begin{enumerate}
    \item \textbf{(Conceptual)} Why does development mode prefix resource names, and what breaks without it?
    \item \textbf{(Conceptual)} Explain the two-command clean-clone redeploy and what each command owns.
    \item \textbf{(Hands-on)} Complete the Thursday lab, including the deliberately-broken PR.
    \item \textbf{(Hands-on)} Add a drift-detection job comparing deployed jobs against the bundle registry.
    \item \textbf{(Hands-on)} Build the per-domain DBU dashboard and write this month's variance note.
    \item \textbf{(Design)} Write the hotfix runbook (SLA, approvers, evidence) for the Friday scenario.
    \item \textbf{(Operations)} Draft the incident review template for ``prod changed outside the pipeline.''
    \item \textbf{(Architecture)} When would you introduce dbt into this delivery model, and which bundle resources would it displace?
\end{enumerate}
